% Project proposal. This is the whole document for now; the full sections
% (01--13) are written later and swapped in from main.tex.

\section{Introduction}\label{sec:intro}

Cryptography needs random numbers. Keys, nonces, session tokens, padding:
all of them have to be numbers an attacker cannot guess. But ``random'' means two
different things, and it is easy to mix them up.
\begin{enumerate}
  \item A sequence \emph{looks} random. It has about as many zeros as ones,
        the runs are the right length, nothing repeats.
  \item A sequence \emph{is} unpredictable. Even knowing everything else, you
        cannot say what the next bit will be.
\end{enumerate}
The first is a question for statistics. The second is a question for
cryptography. A sequence can pass every statistical test and still be
completely predictable. And the other way around: a ciphertext can look like
noise and still leak information about the message inside it.

This project is about that gap. I want to study it with the tools from the
probability course, and I want to check everything by simulation, not just
argue it on paper. So I use small systems I can fully analyse---coin flips, a
linear congruential generator, a Caesar cipher, XOR, the one-time pad---rather
than real ciphers like AES, where nothing can be worked out by hand.

\subsection{What I want to find out}\label{sec:questions}

The main question is:
\begin{quote}
  \textbf{If the randomness used to make an encryption key gets worse, how
  much more does the ciphertext reveal about the message?}
\end{quote}
Two smaller questions come before and after it:
\begin{enumerate}
  \item How do you test, with statistics, whether a sequence behaves
        randomly? And what do those tests miss?
  \item If a sequence passes the tests, is it safe to use as a key?
\end{enumerate}
I expect the answer to the last one is no.  The point of the project is to
show why.

\section{What the project will cover}\label{sec:scope}

There are four parts, plus a small interactive program. Here I only say what
each part does and which piece of probability it rests on. The derivations
and the results go in the full report.

\subsection{Part I: Randomness, and how to test for it}

I model a binary sequence $b_1 b_2 \cdots b_n$ as independent Bernoulli
trials with $\Prob{b_i = 1} = p$. An ideal source has $p = \tfrac12$ and no
dependence between bits. I will generate sequences from several sources and
compare them:
\begin{itemize}
  \item fair and biased coin flips (simulated),
  \item sequences typed by people who were asked to ``be random'',
  \item a linear congruential generator (LCG), which just iterates
        $X_{n+1} = (aX_n + c) \bmod m$,
  \item the operating system's cryptographic generator.
\end{itemize}
Each sequence goes through a small set of tests, simplified from the NIST
suite~\cite{nistsp80022}:
\begin{itemize}
  \item \textbf{Frequency test.} If the bits are fair, the number of ones is
        $X \sim \mathrm{Bin}(n, \tfrac12)$, so $\Ex{X} = n/2$ and
        $\Var{X} = n/4$. The normal approximation gives a $p$-value.
  \item \textbf{Runs test.} Count the runs of identical bits and compare with
        what independence predicts (this is the Wald--Wolfowitz
        idea~\cite{wald1940}).
  \item \textbf{Entropy.} $\Ent{X} = -\sum_x \Prob{x}\log_2\Prob{x}$. It is
        one bit only when the source is fair and independent.
\end{itemize}
The human sequences are the interesting case. People avoid long runs and
switch between 0 and 1 too often, because to them that looks ``more random''.
A two-state Markov chain whose probability of switching is above $\tfrac12$
models this well, and the tests above should catch it.

\subsection{Part II: Encryption, and what it leaks}

I use three ciphers. I chose them because it is easy to see what they do to
the statistics of the message.
\begin{itemize}
  \item The \textbf{Caesar cipher} shifts every letter by the same amount.
        Letter frequencies survive untouched, so frequency analysis breaks
        it.
  \item \textbf{XOR encryption}, $\ctxt = \msg \xor \key$, with keys from
        sources of different quality.
  \item The \textbf{one-time pad}: XOR with a key that is uniformly random,
        independent of the message, as long as the message, and never
        reused.
\end{itemize}
The standard I measure them against is Shannon's definition of perfect
secrecy~\cite{shannon1949}:
\[
  \Prob{\msg = m \mid \ctxt = c} = \Prob{\msg = m}
  \quad\text{for every } m \text{ and } c .
\]
In words: seeing the ciphertext tells you nothing new about the message. This
is a statement about conditional probability, which is where the project
touches the course most directly. I will show that the one-time pad satisfies
it and the other two do not.

Failing perfect secrecy is one thing; I also want a number for \emph{how
badly} a cipher fails. For that I use entropy, conditional entropy and mutual
information, from Shannon's 1948 paper~\cite{shannon1948} in the notation of
Cover and Thomas~\cite{cover2006}:
\[
  \MI{\msg}{\ctxt} = \Ent{\msg} - \CEnt{\msg}{\ctxt} .
\]
Perfect secrecy is the same as $\MI{\msg}{\ctxt} = 0$. Anything above zero is
information that leaked. All of these can be estimated from simulated
$(\msg, \ctxt)$ pairs.

\subsection{Part III: The main experiment---making the key worse}

I take a tiny message space, $\msg \in \bits^2$, and give it a lopsided
distribution on purpose, say
$\Prob{00} = 0.4,\ \Prob{01} = 0.3,\ \Prob{10} = 0.2,\ \Prob{11} = 0.1$.
I draw many messages from it and encrypt each one as $\ctxt = \msg \xor \key$.
The cipher never changes. The only thing I change is where the key comes from.
\begin{description}
  \item[A. Uniform key.] Each of the four keys with probability $\tfrac14$.
        I expect $\MI{\msg}{\ctxt}$ to go to zero as I draw more samples.
  \item[B. Sweep.] Each key bit is $1$ with probability
        $q \in \{0.50, 0.55, \dots, 0.95\}$. For each $q$ I record the key
        entropy, the test results, the ciphertext distribution,
        $\CEnt{\msg}{\ctxt}$ and $\MI{\msg}{\ctxt}$.
\end{description}
Part B gives the main figure of the project: leakage plotted against key
bias, and against key entropy. Because the system is only two bits, I can
also compute the mutual information exactly and check the simulation against
it.

\subsection{Part IV: Looking random is not the same as being secure}

Then I go back to the generators from Part~I. An LCG with well-chosen
constants passes every test in my set. But its whole future is fixed by its
current state, and you can recover that state from a few outputs. So it looks
random and is not secure. Generators built for cryptography are designed so
that recovering the state is computationally infeasible, not just so that
the output looks nice. Knuth's chapter on random numbers~\cite{knuth1997}
and Katz and Lindell's textbook~\cite{katz2020} both make this distinction,
and this part follows them.

\subsection{The interactive part}

Alongside the report I will write a small program (command line or a web
page). It lets you
\begin{itemize}
  \item type in a binary sequence and see all the tests run on it,
  \item try to ``be random'' yourself and get compared to a machine,
  \item look at a few sequences and guess which one was really random,
  \item encrypt a message with each cipher and look at the ciphertext's
        statistics,
  \item play with the leakage-versus-bias curve from Part~III.
\end{itemize}

\section{Probability and statistics I will use}\label{sec:topics}

I want to use the course material, not repeat it.
Table~\ref{tab:topics} lists what the project needs and where Blitzstein and
Hwang~\cite{blitzstein2015}, the course textbook, cover it. Entropy and
mutual information are not in that book; for those I go to
Shannon~\cite{shannon1948} and Cover and Thomas~\cite{cover2006}.

\begin{table}[htbp]
  \centering
  \small
  \begin{tabular}{@{}>{\raggedright\arraybackslash}p{0.30\linewidth}>{\raggedright\arraybackslash}p{0.44\linewidth}>{\raggedright\arraybackslash}p{0.20\linewidth}@{}}
    \toprule
    Concept & Where I use it & Blitzstein \& Hwang \\
    \midrule
    Sample spaces; naive definition of probability
      & Message space, key space, counting the keys that send $m$ to $c$
      & \S1.2--1.3, \S1.6 \\
    Conditional probability; Bayes' rule; law of total probability
      & Definition of perfect secrecy; what you believe about the message
        after seeing the ciphertext; ciphertext distribution as a mixture
        over keys
      & \S2.2--2.3 \\
    Independence of events and of random variables
      & Key independent of message; independent bits in an ideal source
      & \S2.5, \S3.8 \\
    Random variables and PMFs
      & Everything
      & \S3.1--3.2 \\
    Bernoulli and Binomial
      & One bit; number of ones in $n$ bits (frequency test)
      & \S3.3 \\
    Discrete Uniform
      & The one-time-pad key
      & \S3.5 \\
    Expectation; linearity; indicator variables
      & Expected number of ones; expected number of runs (a sum of
        indicators)
      & \S4.1--4.2, \S4.4 \\
    Variance
      & Spread of the frequency and runs statistics when the source is fair
      & \S4.6 \\
    Joint, marginal and conditional distributions
      & The joint law of $(\msg, \key, \ctxt)$; $\CEnt{\msg}{\ctxt}$
      & \S7.1 \\
    Law of large numbers
      & Why estimated frequencies and entropies settle down as $n$ grows
      & \S10.2 \\
    Central limit theorem
      & The normal approximation behind the frequency and runs $p$-values
      & \S10.3 \\
    Markov chains
      & A model of a person who switches between 0 and 1 too often
      & \S11.1 \\
    Simulation
      & Every experiment is a Monte Carlo estimate
      & the ``R'' section of each chapter \\
    \bottomrule
  \end{tabular}
  \caption{Probability concepts the project uses, and where the course
    textbook covers them.}
  \label{tab:topics}
\end{table}

Hypothesis tests and $p$-values run through all of Part~I. The textbook does
not treat them on their own, so the report will state the one definition it
needs---the probability, if the source really were random, of a statistic at
least as extreme as the one I saw---and rely on the distributions above.

\section{Sources}\label{sec:sources}

The project leans on a small number of sources. Here is what I take from
each one.

\paragraph{Shannon, \emph{Communication Theory of Secrecy Systems}
  (1949)~\cite{shannon1949}.}
The definition of perfect secrecy, the proof that the one-time pad has it,
and the result that perfect secrecy needs a key with at least as much entropy
as the message. The whole second half of the project rests on this paper.

\paragraph{Shannon, \emph{A Mathematical Theory of Communication}
  (1948)~\cite{shannon1948}.}
Entropy, conditional entropy and mutual information. These are the numbers I
use to measure leakage.

\paragraph{Cover and Thomas, \emph{Elements of Information Theory}
  (2006)~\cite{cover2006}.}
The modern notation and the basic identities I need in derivations: the chain
rule, $\MI{\msg}{\ctxt} = \Ent{\msg} - \CEnt{\msg}{\ctxt}$, and the fact that
mutual information is never negative.

\paragraph{NIST SP 800-22 (2010)~\cite{nistsp80022}.}
My frequency and runs tests are simplified versions of tests in this suite.
I follow its conventions for $p$-values.

\paragraph{Knuth, \emph{The Art of Computer Programming}, Vol.~2, Chapter~3
  (1997)~\cite{knuth1997}.}
Linear congruential generators, how to choose their constants, and the
classical empirical tests for randomness.

\paragraph{Katz and Lindell, \emph{Introduction to Modern Cryptography}
  (2020)~\cite{katz2020}.}
The textbook treatment of perfect secrecy, and the argument for why ``passes
the statistical tests'' is not a definition of security.

\paragraph{Wald and Wolfowitz (1940)~\cite{wald1940}.}
The original paper for the runs test.

\paragraph{Blitzstein and Hwang, \emph{Introduction to Probability}
  (2015)~\cite{blitzstein2015}.}
The course textbook. All the probability background is referred to it, as in
Table~\ref{tab:topics}.

\section{A first derivation}\label{sec:first-derivation}

Here is the one derivation the whole project rests on.  I do it now to
show that the notation above works.

Take $\msg, \key, \ctxt \in \bits^n$ with $\ctxt = \msg \xor \key$, and let
$\key$ be uniform on $\bits^n$ and independent of $\msg$. Fix any $m$ and $c$.
Because XOR is its own inverse, the event $\{\msg = m, \ctxt = c\}$ is the
same as $\{\msg = m, \key = m \xor c\}$, so by independence
\[
  \Prob{\ctxt = c \mid \msg = m} = \Prob{\key = m \xor c} = 2^{-n}.
\]
The right-hand side does not depend on $m$. By the law of total probability,
\[
  \Prob{\ctxt = c} = \sum_{m'} \Prob{\ctxt = c \mid \msg = m'}\Prob{\msg = m'}
                   = 2^{-n} \sum_{m'} \Prob{\msg = m'} = 2^{-n},
\]
and Bayes' rule finishes it:
\[
  \Prob{\msg = m \mid \ctxt = c}
    = \frac{\Prob{\ctxt = c \mid \msg = m}\,\Prob{\msg = m}}{\Prob{\ctxt = c}}
    = \frac{2^{-n}\,\Prob{\msg = m}}{2^{-n}} = \Prob{\msg = m}.
\]
That is Shannon's condition.  It holds for every message distribution,
including the lopsided one in Part~III.  The proof uses three tools:
conditional probability, the law of total probability and Bayes' rule.
All three are in Chapter~2 of the course textbook.

The same computation shows what breaks when the key is not uniform.  If
$\Prob{\key = k}$ depends on $k$, then
$\Prob{\ctxt = c \mid \msg = m} = \Prob{\key = m \xor c}$ depends on $m$.
So the ciphertext carries information about the message, and
$\Prob{\msg = m \mid \ctxt = c} \neq \Prob{\msg = m}$ in general.
Part~III measures how much, using $\MI{\msg}{\ctxt}$.

\section{Plan of work}\label{sec:plan}

Table~\ref{tab:plan} is the order I will do things in.  In each part the
simulations come before the writing.  That way the report describes
results I have, not results I expect.

\begin{table}[htbp]
  \centering
  \small
  \begin{tabular}{@{}lp{0.62\linewidth}@{}}
    \toprule
    Week & Work \\
    \midrule
    1 & Set up the code: bit-sequence sources (coins, LCG, OS generator)
        and the three tests. Check each test on a known-fair source
        by looking at the distribution of its $p$-values. \\
    2 & Collect human sequences; fit the two-state Markov model; write
        Part~I. \\
    3 & Caesar, XOR and one-time-pad implementations; estimators for
        $\Ent{\cdot}$, $\CEnt{\cdot}{\cdot}$ and $\MI{\cdot}{\cdot}$ from
        samples; write Part~II. \\
    4 & Main experiment (A and the sweep B); exact mutual information
        for the two-bit system to compare against; the main figure. \\
    5 & Part~IV: recover an LCG's state from its outputs. Interactive
        program. \\
    6 & Write the introduction and conclusion last; proofread; final
        figures. \\
    \bottomrule
  \end{tabular}
  \caption{Planned order of work.}
  \label{tab:plan}
\end{table}

\section{The end goal}\label{sec:goal}

I want to follow the chain
\[
  \text{randomness} \;\longrightarrow\; \text{uncertainty}
  \;\longrightarrow\; \text{information} \;\longrightarrow\; \text{secrecy}
\]
from one end to the other, with each link backed by both a derivation and a
simulation. Concretely, I expect to show that:
\begin{enumerate}
  \item sequences people make up by hand are detectably not random;
  \item simple statistical tests catch biased and badly built generators;
  \item a Caesar cipher keeps the letter statistics of the plaintext, so its
        ciphertext leaks nearly all of $\Ent{\msg}$;
  \item the one-time pad with a uniform key gives $\MI{\msg}{\ctxt} \approx 0$
        in simulation, matching the exact answer of zero;
  \item biasing the key makes $\MI{\msg}{\ctxt}$ positive, and the leakage
        keeps growing as the key's entropy falls;
  \item an LCG can pass every test in my set and still be predictable from a
        handful of outputs.
\end{enumerate}
Put together, these support the claim the project is named after:
\begin{quote}
  \textbf{Good randomness is necessary for secrecy. But looking random is not
  enough to make something secure.}
\end{quote}
